\chapter{Nonlinear Power Excursions}
\label{ch:nonlinear}

Linear analysis tells us whether an operating point is stable. When a reactor is
driven hard---a large, fast reactivity insertion---the response is fully
nonlinear, and the question becomes: how high does the power spike, how much
energy is deposited, and what stops it? The classic answer is the
Nordheim--Fuchs model, which gives a remarkably clean closed-form description of a
self-limiting prompt excursion. This is the model that quantifies what a
\emph{thermally stable} reactor does when subjected to a reactivity-initiated
accident---and, by its failure, what Chernobyl did.

\section{The reactivity-initiated accident}

A reactivity-initiated accident (RIA) is a rapid insertion of positive
reactivity---an ejected control rod, a cold-water slug, or, at Chernobyl, the
catastrophic positive scram. If the insertion exceeds one dollar, the reactor
becomes prompt-critical and the power rises on the $\Lambda$ timescale, far too
fast for delayed neutrons or mechanical systems to matter. The only thing that can
arrest it is prompt feedback. Whether the excursion is a benign, self-terminating
spike or an explosive energy release depends entirely on the sign and strength of
that prompt feedback.

\section{The Nordheim--Fuchs model}

Consider a step insertion $\rho_0>\beta$ on a prompt-critical reactor with prompt
negative feedback. On the prompt timescale the delayed neutrons are a nearly
constant background, so we drop the precursor term and write the prompt kinetics
\begin{equation}
\frac{\dd p}{\dd t} = \frac{\rho(t)-\beta}{\Lambda}\,p,
\end{equation}
with feedback through an adiabatically heating fuel (Chapter~\ref{ch:thermal}),
\begin{equation}
\rho(t) = \rho_0 - \gamma\,\big(E(t)\big),\qquad
\frac{\dd E}{\dd t} = p,
\end{equation}
where $E(t)=\int_0^t p\,\dd t'$ is the energy released and $\gamma>0$ is the
magnitude of the (negative) feedback per unit energy---the product of the fuel
heat capacity and the temperature coefficient, $\gamma=|\alpha|/C$. Writing the
prompt reactivity above critical as $\rho_0-\beta$ and feedback $-\gamma E$, the
power is
\begin{equation}
\frac{\dd p}{\dd t} = \frac{(\rho_0-\beta) - \gamma E}{\Lambda}\,p.
\end{equation}
Using $\dd p/\dd t = (\dd p/\dd E)\,p$ and integrating once,
\begin{equation}
p(E) = p_0 + \frac{1}{\Lambda}\!\left[(\rho_0-\beta)E - \tfrac12\gamma E^2\right].
\end{equation}

\section{Peak power, energy, and temperature rise}

The power peaks when $\dd p/\dd t=0$, i.e.\ when the feedback exactly cancels the
prompt reactivity, $\gamma E_{\text{peak}} = \rho_0-\beta$, so the energy released
at the peak is
\begin{equation}
\boxed{\ E_{\text{peak}} = \frac{\rho_0-\beta}{\gamma}\ }
\end{equation}
The excursion does not stop here---the power is maximal, not zero---but the
feedback now drives the reactivity below prompt-critical and the power falls. The
total energy released as the prompt spike completes is, by the symmetry of the
parabola $p(E)$, twice the energy at the peak:
\begin{equation}
\boxed{\ E_{\text{total}} \approx \frac{2(\rho_0-\beta)}{\gamma}\ }
\end{equation}
and the peak power is
\begin{equation}
p_{\max} \approx \frac{(\rho_0-\beta)^2}{2\gamma\Lambda}.
\end{equation}
The corresponding adiabatic fuel temperature rise is
$\Delta \Tf = E_{\text{total}}/C = 2(\rho_0-\beta)/(\gamma C)$. These formulae are
the quantitative essence of self-limitation: the stronger the negative feedback
$\gamma$, the smaller the energy release and temperature rise; the larger the
prompt reactivity excess $\rho_0-\beta$, the larger and sharper the spike.

\begin{example}[A benign spike vs.\ a destructive one]
With prompt feedback strong enough that $\gamma$ limits $\Delta \Tf$ to a few
hundred kelvin, an RIA produces a fast spike that the fuel survives---this is the
designed-for behaviour validated in tests such as SPERT and BORAX. If instead the
\emph{net} prompt feedback is weak, near zero, or---catastrophically---positive,
$E_{\text{total}}$ grows without the parabolic turnover, the temperature rise is
unbounded in the model, and in reality the energy deposition fragments the fuel,
flashes the coolant, and disassembles the core. Chernobyl's excursion deposited
enough energy in a fraction of a second to shatter the fuel and trigger a steam
explosion, precisely because the prompt feedback in its configuration was not
saving-ly negative.
\end{example}

\section{Self-disassembly and the ultimate shutdown}

In the most violent excursions a second, brutal feedback eventually intervenes:
the core physically disassembles. Thermal expansion, fuel vaporisation, and
pressure-driven dispersal increase neutron leakage and separate the fissile
material, terminating the chain reaction---the Fuchs--Hansen ``disassembly''
mechanism that ended the historical criticality excursions (SL-1, the Godiva and
SPERT destructive tests). Disassembly is a shutdown mechanism only in the sense
that an explosion ends a chain reaction: it does so by destroying the reactor. The
purpose of thermal stability and prompt negative feedback is to terminate the
excursion by gentle means---temperature feedback---long before disassembly is
reached.

\section{Connection to the linear picture}

The Nordheim--Fuchs excursion and the linear stability of Chapter~\ref{ch:linear}
are two limits of the same coupled system. Linear analysis governs small
perturbations about an operating point and asks whether they grow; the
Nordheim--Fuchs analysis governs the large, fast, prompt-critical transient and
asks how the nonlinear feedback caps it. A reactor that is linearly stable with
strong negative feedback is also the one whose Nordheim--Fuchs excursions are
mild. A reactor that is linearly unstable---positive feedback---has no benign
nonlinear limit; its excursions terminate only by disassembly. The companion
computations of Chapter~\ref{ch:comp} show both behaviours side by side.

\keypoint{In a prompt-critical excursion, prompt negative feedback caps the energy
release at $E_{\text{total}}\approx 2(\rho_0-\beta)/\gamma$ and the reactor
survives. Without net negative prompt feedback there is no such cap, and the
excursion ends only by destroying the core. Thermal stability is the difference
between a spike and an explosion.}


\section*{Exercises}
\addcontentsline{toc}{section}{Exercises}
\begin{enumerate}[leftmargin=*,itemsep=2pt]
\item Derive $E_{\text{peak}}=(\rho_0-\beta)/\gamma$ from the Nordheim--Fuchs model.
\item Show that the total energy release is approximately $2(\rho_0-\beta)/\gamma$.
\item How does the peak power scale with the prompt reactivity excess and the feedback strength?
\item What physically terminates an excursion when prompt feedback is absent?
\item Relate the sign of the feedback in the Nordheim--Fuchs model to the linear stability of Chapter 8.
\end{enumerate}
